% SEGMENT OLP-0611-S01
% Part: methods
% Chapter: proofs
% Section: cant-do-it

\documentclass[../../../include/open-logic-section]{subfiles}

\begin{document}

\olfileid{mth}{prf}{cnt}
\olsection{என்னால் செய்ய முடியவில்லை!{}}

கைவிட்டுவிடலாம் என்று தோன்றும் கட்டத்தை நாம்
எல்லோரும் சந்திக்கிறோம். ஆனால் உங்களால் இதைச்
\emph{செய்ய முடியும்}. உங்கள் ஆசிரியரும் கற்பித்தல்
உதவியாளரும் உடன் படிப்பவர்களும் உதவ முடியும்.
அவர்களிடம் உதவி கேளுங்கள்!{} நெருக்கடியைத் தவிர்க்கவும்,
கைவிடத் தோன்றும்போது சமாளிக்கவும் சில ஆலோசனைகள் இதோ.

கணக்குகளை வெற்றிகரமாகத் தீர்க்கப் பின்வருவனவற்றைச் செய்க:
\begin{enumerate}
% SEGMENT OLP-0611-S02
\item \emph{முடிந்தவரை முன்கூட்டியே தொடங்குக.}
  பருவம் முழுவதும் வேலைகள் குவியும்; பலருக்கும்
  தள்ளிப்போடும் பழக்கத்தை வெல்வது கடினம். வீட்டுப்
  பயிற்சிகளை முன்கூட்டியே தொடங்குவது உங்களால் செய்ய
  முடிந்த சிறந்த காரியங்களில் ஒன்று. அப்போது ஒரு
  கணக்கில் சிக்கினாலும், அழுதுகொண்டிருப்பதைத் தவிர
  தீர்வு தேட நேரம் இருக்கும்.
% SEGMENT OLP-0611-S03
\item \emph{உடன் படிப்பவர்களிடம் பேசுக}. நீங்கள்
  தனியாக இல்லை. வகுப்பில் மற்றவர்களுக்கும் சிரமம்
  இருக்கலாம்; அவர்களுக்கு வரும் சிரமம் வேறானதாக
  இருக்கலாம். அவர்களுடன் பேசுவது கணக்கை வேறு
  கோணத்தில் பார்க்க உதவும்; அதனால் முன்னேற்றம்
  கிடைக்கலாம். அவர்களுடைய தீர்வை அப்படியே
  எழுதிவிடாதீர்கள். ஒரு குறிப்பைக் கேளுங்கள்;
  அல்லது எங்கே சிக்குகிறீர்கள் என்று விளக்கி
  அடுத்த படி என்னவென்று கேளுங்கள். பிறகு உங்களால்
  முடிந்தபோது அவர்களுக்கும் உதவுங்கள். மற்றொருவருக்கு
  விளக்கிச் சொல்வது உங்களுடைய புரிதலையும் வளர்க்கும்.
% SEGMENT OLP-0611-S04
\item \emph{உதவி கேளுங்கள்.} உங்களுக்கு உதவும்
  பல வழிகள் உள்ளன. ஆசிரியரும் கற்பித்தல் உதவியாளரும்
  உங்களுக்காகவே இருக்கிறார்கள்; நீங்கள் \emph{வெற்றி
  பெற வேண்டும்} என்பதே அவர்களின் விருப்பம். ஒரு
  கணக்கை எப்படிச் செய்வது என்றும் எந்தப் படியில்
  சிக்கல் ஏற்படுகிறது என்றும் கண்டறிய அவர்கள் உதவலாம்.
% SEGMENT OLP-0611-S05
\item \emph{சிறிது ஓய்வெடுங்கள்.} சிக்கியதற்குக்
  காரணம் ஒரே கணக்கை மிக நீண்ட நேரம் பார்த்துக்
  கொண்டிருப்பதாகவும் \emph{இருக்கலாம்}. சிறிது
  இடைவெளி எடுங்கள்; தேநீர் அருந்துங்கள்; வேறொரு
  கணக்கைச் சில நேரம் செய்து பாருங்கள். பிறகு
  புத்துணர்ச்சியுடன் இதற்குத் திரும்புங்கள். ஒரு
  இரவு தூங்கி எழுந்தும் யோசித்துப் பாருங்கள்.
\end{enumerate}

% SEGMENT OLP-0611-S06
இந்த வழிகள் எல்லாம் நிறுவலைப் பற்றி முன்கூட்டியே
வேலை செய்யத் தொடங்கியிருந்தால் மட்டுமே பயனளிக்கும்
என்பதைக் கவனித்தீர்களா? ஒப்படைக்க வேண்டிய நாளுக்கு
முன் அதிகாலை இரண்டு மணிக்குத்தான் நிறுவலைத்
தொடங்கினால், இவை பெரிதும் உதவாமல் போகலாம்.

இது அச்சமூட்டுவது போலத் தோன்றலாம். ஆனால் ஒரு
நிறுவலைத் தீர்ப்பது இறுதியில் பலனளிக்கும் சவால்.
சிலர் இதையே தொழிலாகச் செய்கிறார்கள்; இதில்
மகிழ்ச்சியளிக்கும் ஏதோ ஒன்று இருக்கவேண்டும்.
மற்ற எல்லாத் திறன்களைப் போலவே கணக்குகளைத்
தீர்ப்பதற்கும் நிறுவல்களை எழுதுவதற்கும் பயிற்சி
தேவை. இதை எளிதாகச் செய்யும் சக மாணவர்களைப்
பார்க்கலாம்; அவர்கள் ஏற்கெனவே அதிகம் பயிற்சி
செய்திருக்கக்கூடும். சீக்கிரம் மனம் தளராதீர்கள்.

% SEGMENT OLP-0611-S07
ஒரு குறிப்பிட்ட கணக்கிற்கான நேரமோ பொறுமையோ
தீர்ந்துவிட்டாலும் பரவாயில்லை. அதனால் நீங்கள்
அறிவற்றவர் என்றோ, அதை ஒருபோதும் புரிந்துகொள்ள
மாட்டீர்கள் என்றோ பொருள் இல்லை. ஆசிரியரிடமோ
மற்றொரு மாணவரிடமோ அதை எப்படிச் செய்வது என்று
கேட்டறியுங்கள். எங்கே தவறு செய்தீர்கள் அல்லது
சிக்கினீர்கள் என்று கண்டுபிடியுங்கள்; அடுத்த முறை
அதேபோன்ற சிக்கலைத் தவிர்க்க அது உதவும். பின்னர்
தீர்வைப் பார்க்காமல் நீங்களே செய்து பாருங்கள்.
அடுத்த முறை முன்கூட்டியே தொடங்கி, உதவியையும்
முன்கூட்டியே கேளுங்கள்.

\end{document}
